\scanpage{28}
\textbf{C. R. Acad. Sc. Paris, t. 279 (22 juillet 1974)}\hfill\textbf{Série A ---}\mspage{131}\par\noindent\rule{\linewidth}{.35pt}\par

et \((\gamma,\omega)\) un élément de \(\eta\), les \(\eta\)-boucles de \(\mathrm{K}\) sont d'une part les boucles dont on obtient un circuit représentatif par la loi qui étant donné un arc \(a\) et \(i\in\mathbf{Z}_2\) leur associe l'arc \((-1)^j\,\omega_s(-a)\) et \(j\in\mathbf{Z}_2\), où \(s\) est l'extrémité de \(a\) et \(j=i+\gamma(a)\), et d'autre part les circuits réduits à un arc boucle isolé \(b\) telle que \(\gamma(b)=0\). Cette définition ne dépend pas du représentant \((\gamma,\omega)\) choisi. On note \(\mathrm{C}(\mathrm{K},\eta)\) l'ensemble des \(\eta\)-boucles et on a une bijection canonique \(\psi\) entre \(\mathrm{C}(\mathrm{K},\eta)\) et \(\pi_0\,\partial\mathrm{T}(\mathrm{K},\eta)\).

12. \textsc{Classification des Plongements.}

12.1. \textsc{Définition.} --- Deux plongements \(f_1\) et \(f_2\) de \(\mathrm{K}\) dans \(\mathrm{V}_1\) et \(\mathrm{V}_2\) (deux surfaces compactes) sont \textit{équivalents} s'il existe un homéomorphisme \(h\) de \(\mathrm{V}_1\) sur \(\mathrm{V}_2\) tel que \(h\circ f_1=f_2\).

12.2. Soit un plongement \(f\) de \(\mathrm{K}\) dans la surface compacte \(\mathrm{V}\) tel que \(\pi_0(f)\) soit surjectif. Il lui est associé un premier invariant (à équivalence près) la cocircularisation \(\eta_f\) induite sur \(\mathrm{K}\) par un voisinage tubulaire \(\mathrm{T}\) quelconque de \(f(\mathrm{K})\) dans \(\mathrm{V}\).

12.3. Soit alors \(\Sigma=\mathrm{V}\backslash\mathrm{int}\,\mathrm{T}\) ; \(\Sigma\) est une surface à bord \(\partial\Sigma=\partial\mathrm{T}\). On obtient une partition de \(\pi_0\,\partial\Sigma\) en rangeant dans une même classe \(\mathrm{E}_i\) les composantes de \(\partial\Sigma\) contenues dans une même composante \(\Sigma_i\) de \(\Sigma\). Ceci induit une partition de \(\mathrm{C}(\mathrm{K},\eta_f)\) qui permet de construire un second invariant du plongement. A la classe \(\mathrm{E}_i\) est associé un entier \(+1\) ou \(-1\) selon que \(\Sigma_i\) est, ou non, orientable et un autre entier : le genre de \(\Sigma_i\). Tous ces invariants permettent la classification (un peu technique) des plongements compte tenu des « symétries » du voisinage tubulaire, du fait que si \(\Sigma_i\) est orientable, tout homéomorphisme de \(\partial\Sigma_i\) ne se prolonge pas forcément à \(\Sigma_i\) et que, par contre :

12.4. \textsc{Proposition.} --- \textit{Si \(\mathrm{W}\) est une surface \textup{non orientable} compacte à bord \(\partial\mathrm{W}\) et \(h\) un homéomorphisme de \(\partial\mathrm{W}\), alors \(h\) s'étend en un homéomorphisme de \(\mathrm{W}'\).}

13. \textsc{Plongements minimaux.}

13.1. Avec les notations de 12 ci-dessus, si \(\chi\) désigne la caractéristique d'Euler-Poincaré on a \(\chi(\mathrm{V})=\chi(\mathrm{K})+\chi(\Sigma)\).

13.2. Un plongement \(f\) de \(\mathrm{K}\) dans une surface orientable connexe compacte \(\mathrm{V}\) est \textit{minimal} si pour tout plongement dans une autre surface \(\mathrm{V}'\) le genre de \(\mathrm{V}'\) est supérieur ou égal à celui de \(\mathrm{V}\). Le genre de \(\mathrm{K}\) est alors, par définition, celui de \(\mathrm{V}\).

13.3. On note \(\Omega(\mathrm{K})=\sup\limits_\omega\mathrm{Card}\,\mathrm{C}(\mathrm{K},\omega)\) [\textit{cf.} (\textsuperscript{1})], c'est un entier positif attaché à \(\mathrm{K}\) et on a :

\textsc{Proposition.} --- \textit{Si \(\mathrm{K}\) est connexe son genre est \(1-(1/2)\,(\chi(\mathrm{K})+\Omega(\mathrm{K}))\).}

Cela conduit à l'algorithme d'Edmonds (\textsuperscript{2}).

13.4. Si \(\mathrm{K}\) a pour composantes connexes \(\mathrm{K}_1,\mathrm{K}_2,..,\mathrm{K}_p\) le genre de \(\mathrm{K}\) est la somme des genres de \(\mathrm{K}_1,\mathrm{K}_2,\ldots,\mathrm{K}_p\). On retrouve un résultat connu (\textsuperscript{3}) ; de plus :

13.5. \textsc{Théorème.} --- \textit{Avec les notations précédentes \textup{(13.4)} les plongements minimaux de \(\mathrm{K}\) sont ceux obtenus de la façon suivante : soient \(f_i\) (\(i=1,\ldots,p\)) des plongements minimaux des \(\mathrm{K}_i\) dans des surfaces compactes \(\mathrm{V}_i\), et soit \(\Gamma\) une structure d'arbre sur l'ensemble des \(\mathrm{V}_i\), on forme une surface \(\mathrm{V}\) en « reliant par un cylindre » tout couple \((\mathrm{V}_i,\mathrm{V}_j)\) tel que \(\mathrm{V}_i\) soit lié à \(\mathrm{V}_j\) dans \(\Gamma\).}

{\small (*) Séance du 5 juin 1974.}
% archive sheet number: 27

